\chapter{Search Methods}
\label{chap:algorithms}

% Sources: src/hrgs_scheduler/search/, docs/Optimizer Status.md,
% docs/Justification of Implementation.md §4

Three search tiers share one schedule representation and one evaluator. Whichever tier is under discussion, the schedule it returns is called \emph{the optimizer} throughout this thesis; no single software artefact carries that name, and Chapters~\ref{chap:results}--\ref{chap:discussion} state which tier produced a given result whenever it matters. Every returned candidate passes the structural checks of Section~\ref{sec:model:legality} and is evaluated by the bottom-up pass of Section~\ref{sec:model:objectives}; the tiers differ only in which legal \Glspl{dag} they consider, never in how a returned \Gls{dag}'s fidelity, success probability, latency, or cost is calculated.

The chapter moves from general search concepts (Section~\ref{sec:alg:prelims}) through the shared evaluator (Section~\ref{sec:alg:evaluator}) to three increasingly broad but less exhaustive tiers: fixed baseline families (Section~\ref{sec:alg:bruteforce}), the exact span-partition recurrence (Section~\ref{sec:alg:dp}), same-span pumping (Section~\ref{sec:alg:pumping}), and beam search (Section~\ref{sec:alg:beam}), the heuristic used for the production-scale experiments of Chapter~\ref{chap:results}.

\section{Preliminaries: Combinatorial Search Concepts}
\label{sec:alg:prelims}

The schedule space defined in Section~\ref{sec:model:dag} is combinatorially large: even modest chains admit too many legal schedules to enumerate individually. The three tiers below rely on three standard tools from combinatorial optimisation, introduced briefly here before being applied to schedules.

\paragraph{Multi-objective optimisation and Pareto dominance} A search problem is \emph{multi-objective} when candidates are compared on several quantities at once, here cost $C$, fidelity $F$, and success probability $P_{\mathrm{success}}$ (Section \ref{sec:model:objectives}). Candidate $A$ \emph{Pareto-dominates} $B$ if $A$ is at least as good on every quantity and strictly better on one; a \emph{Pareto-optimal} candidate is dominated by none, and the \emph{Pareto frontier} is the set of all such candidates~\cite{Pareto2014Manual,Deb2002Fast,Ehrgott2005Multicriteria}. Retaining the whole frontier matters here because a locally inferior candidate can become the better choice once combined with another resource, as happens when two spans are swapped together.

\paragraph{Dynamic programming} \emph{Dynamic programming} decomposes a large problem into smaller overlapping subproblems that are solved once and reused wherever they recur (\emph{memoized})~\cite{Bellman1984Dynamic}. Section~\ref{sec:alg:dp} applies this principle to schedules: candidate sub-schedules for a short span are computed once and reused when constructing longer spans.

\paragraph{Heuristic and beam search} When candidates become too numerous to enumerate exhaustively, a \emph{heuristic search} trades completeness for tractability by discarding some before they are fully explored. \emph{Beam search} does this by keeping only a fixed number of candidates, the \emph{beam width}, at each stage, ranked by a chosen rule~\cite{Lowerre1976HARPY,Huang2018When,Meister2022BestFirst}; a narrower beam is cheaper but more likely to discard a candidate that would have led to a better final schedule. Section~\ref{sec:alg:beam} applies this to the span-partition recursion of Section~\ref{sec:alg:dp}.

\section{The Schedule Evaluator}
\label{sec:alg:evaluator}

% Source: schedule/evaluator.py, experiments/fig5_fidelity_vs_noise.py

The evaluator makes one bottom-up pass over a schedule \Gls{dag}. At each node it applies the corresponding physical operation to the current Pauli-channel state, synchronizes branches that become ready at different times, applies the resulting memory decoherence, and propagates side effects and timing. At the root it returns fidelity $F$, success probability $P_{\mathrm{success}}$, resource cost $C$, and latency $L$. Rate is then $R=P_{\mathrm{success}}/L$ under the full-restart model of Section~\ref{sec:model:objectives}.

Section~\ref{sec:exp:validation} shows that this pass reproduces the source paper's Fig.\,5 fidelity numbers to four significant figures. This validates the physical model implemented by the evaluator, but does not establish that any search tier finds an optimal schedule, which is a separate question addressed by this chapter and Chapter~\ref{chap:results}.

\section{Fixed Schedule Families}
\label{sec:alg:bruteforce}

The first tier enumerates four fixed structural families:
\begin{itemize}
    \item a \emph{single raw end-to-end chain}, with no purification;
    \item \emph{end-node pumping}, in which $n_{\mathrm{pur}}$ independent raw chains are purified sequentially at the end nodes with a round-trip herald between rounds (\emph{heralded} end-node pumping);
    \item its \emph{optimistic counterpart}, deferring heralding until the end (\emph{optimistic} end-node pumping);
    \item \emph{uniform link-level purification}, using the same copy count and circuit at every hop.
\end{itemize}
For even $N$ within budget, the Integrating paper's hand-chosen schedule~\cite{Benchasattabuse2025Integrating} is included as a fifth fixed comparison point.

Three of these families, the raw chain, heralded end-node pumping, and optimistic end-node pumping, are referred to throughout this thesis as the \emph{canonical schedules}: they are the schedules whose fidelity is checked directly against the source paper's published Fig.\,5 benchmarks (Section~\ref{sec:exp:validation}).

Within each family, purification sequences use the three two-copy circuits $\mathcal C_{\mathrm{YY}}$, $\mathcal C_{\mathrm{ZX}}$, and $\mathcal C_{\mathrm{XZ}}$ (Section~\ref{sec:bg:circuits}). All $3^r$ circuit sequences are enumerated up to a configured round cutoff $r_{\max}$; above that cutoff, the implementation uses a small curated set. Thus, ``brute force'' denotes exhaustive enumeration within these predefined structural families and circuit grid, not enumeration of every legal schedule \Gls{dag}.

These fixed families serve two purposes: they provide explicit baselines throughout Chapter~\ref{chap:results} (the raw chain and source-paper schedule as reference points, uniform link-level purification as a simple non-adaptive baseline requiring no structural search), and their candidates are always added to those produced by the two broader search tiers, so a known schedule class is never lost solely to frontier pruning.

\section{Span Dynamic Programming}
\label{sec:alg:dp}

For every span $\mathrm{Span}(a,b)$, the second tier memoizes a Pareto frontier~\cite{Pareto2014Manual} of ways to create one resource over that span. A single-hop frontier ($b=a+1$) contains the raw link and link-level purification choices. For a wider span, the recurrence tries every split point $a<m<b$ and combines each retained candidate from $\mathrm{Span}(a,m)$ with each from $\mathrm{Span}(m,b)$ using a Swap operation (Section~\ref{sec:bg:hrgs}). Each sub-span is computed once and reused in wider constructions. This reuse, rather than repeatedly constructing and evaluating complete end-to-end schedules, is the main computational advantage of the dynamic program.

\clearpage

The frontier retained at each span is multi-objective over $(C,F,P_{\mathrm{success}})$ (Pareto dominance, Section~\ref{sec:alg:prelims}). Only non-dominated candidates are retained, following the standard Pareto-pruning rule~\cite{Ehrgott2005Multicriteria}. This pruning is exact with respect to the three quantities tracked by the frontier, but does not make the search exhaustive: the caps and heuristic retention rules introduced in the next two sections can still discard a non-dominated candidate before it reaches the end-to-end span.

\section{Same-Span Entanglement Pumping}
\label{sec:alg:pumping}

The span-partition recurrence only builds a resource by swapping two \emph{different} spans, so it cannot represent purifying a resource against a fresh, independently generated candidate at the \emph{same} span $\kappa$. The second-tier extension adds this operation explicitly, bounded by configured copy-count and circuit limits rather than left unrestricted. An uncapped mode, restoring exhaustive pumping within those same limits, is available for targeted validation but can already become impractical at small $N$ once the budget is large, so it is used only as a diagnostic, not the general production setting.

Pumping and ordinary Swap candidates compete for the same fixed-size frontier, so adding pumping can evict a swap-only candidate before it reaches an end-to-end frontier. Disabling pumping entirely reproduces the earlier pumping-free search behaviour; Section~\ref{sec:exp:headline} reports both settings at the reference configuration since they can return different top-ranked schedules.

\section{Beam Search}
\label{sec:alg:beam}

The third tier, beam search, reuses the same span recursion, node construction, and evaluator as Section~\ref{sec:alg:dp}, but caps every retained span frontier at a beam width $w_{\text{beam}}$ (default $25$; Appendix~\ref{app:sweeps} justifies this choice), bounding runtime and memory up to $N=28$ hops at the cost of possibly pruning a candidate a wider beam would keep. The cap is necessary in practice: an uncapped width-7 frontier already holds roughly $1{,}200$ candidates, making an uncapped search through the paper's $N=10$ chain impractical. The same $w_{\text{beam}}$ also caps the same-span pumping pool, so pumping and ordinary swap candidates share one set of slots at each span.

At each span, the $w_{\text{beam}}$ slots split between two rankings: roughly half keep the highest-fidelity candidates, preserving purified building blocks that only pay off after a later swap, while the rest favour candidates meeting the fidelity hint $F_{\min}$, then ranked by success probability, fidelity, and cost. This split matters because a purely efficiency-based ranking would reject link-level purification at short spans, where a raw link is locally cheaper -- yet accumulated noise across several such raw links can later make the end-to-end floor unattainable, which the fidelity-first half guards against.

Throughout this thesis, a \emph{score} is the returned candidate's value of whichever quantity the active objective maximizes (Section~\ref{sec:model:objectives}). Under the default rate-maximizing objective used for the beam-search benchmarks of Table~\ref{tab:benchmarks} in Appendix~\ref{app:sweeps}, score is simply the candidate's rate $R(\Sigma)$.

A different heuristic (e.g.\ simulated annealing or a genetic search) would need its own schedule encoding and proposal operators respecting the legality conditions of Section~\ref{sec:model:legality}. Beam search instead keeps the same span-partition representation and only changes how many candidates survive at each frontier, keeping every tier on one physical evaluation path and every run \emph{deterministic}. The result is still only a valid candidate found by a bounded search, not a guarantee of global optimality.